\documentclass{article}
\usepackage{neurips_2025}
\usepackage[utf8]{inputenc}
\usepackage[T1]{fontenc}
\usepackage{amsmath,amssymb,amsthm}
\usepackage{booktabs}
\usepackage{graphicx}
\usepackage{hyperref}
\usepackage{xcolor}
\usepackage{algorithm}
\usepackage{algpseudocode}
\usepackage{nicefrac}

\theoremstyle{plain}
\newtheorem{theorem}{Theorem}
\newtheorem{lemma}{Lemma}
\newtheorem{proposition}{Proposition}

\theoremstyle{definition}
\newtheorem{definition}{Definition}
\newtheorem{assumption}{Assumption}

\title{AIT-LCD: Adaptive Information-Theoretic Local Causal Discovery\\with Explicit Conditioning Set Awareness}

\author{}

\begin{document}

\maketitle

\begin{abstract}
Local causal discovery aims to identify the direct causes and effects of a target variable without learning the complete causal graph. Existing constraint-based methods suffer from unreliable conditional independence (CI) tests in small samples due to fixed significance thresholds, while score-based methods lack computational efficiency for high-dimensional data. We propose AIT-LCD (Adaptive Information-Theoretic Local Causal Discovery), which extends the adaptive thresholding foundations established by \citet{lerner2013adaptive} with explicit conditioning set awareness through a principled functional form $\tau(n,k) = \alpha \cdot \sqrt{k/n} \cdot \log(1 + n/(k \cdot \beta))$. AIT-LCD combines bias-corrected mutual information estimates with adaptive thresholds that explicitly account for both sample size and conditioning set dimensionality. In pilot experiments on the Asia benchmark network, AIT-LCD shows a trend toward improved Markov blanket recovery accuracy (PC F1 of 0.592 vs.~0.548 for IAMB), though this improvement was not statistically significant ($p = 0.078$) given limited statistical power. This workshop paper presents the methodological framework and initial findings from a pilot study, with broader evaluation and ablation study refinement identified as critical next steps.
\end{abstract}

\section{Introduction}\label{sec:intro}

\subsection{Problem Context}

Causal discovery from observational data is a fundamental challenge in machine learning with applications in healthcare, economics, and scientific discovery. While global causal discovery learns the entire causal graph, many real-world scenarios only require understanding the local causal structure around specific target variables. For example, in healthcare, identifying direct risk factors for a specific disease; in genomics, finding direct regulators of a target gene; and in economics, discovering immediate causes of an economic outcome.

Local causal discovery focuses on identifying the Markov blanket (MB) of a target variable---the minimal set of variables that renders the target independent of all others. The MB consists of parents (direct causes), children (direct effects), and spouses (other parents of the children). Learning the MB is sufficient for optimal prediction and provides insights into the local causal structure around the target.

\subsection{Current Limitations}

Existing local causal discovery methods fall into two categories, each with significant drawbacks:

\textbf{Constraint-based methods} (e.g., IAMB \citep{tsamardinos2003algorithms}, HITON-MB \citep{aliferis2003hiton}, PCMB \citep{pena2007towards}) use conditional independence tests to build the Markov blanket iteratively. These methods require exponentially many samples as the conditioning set size grows, suffer from cascading errors when CI tests fail in small samples, and use fixed significance thresholds regardless of sample size.

\textbf{Score-based methods} (e.g., MBMML \citep{tabandeh2021markov}) employ scoring functions to evaluate candidate Markov blankets. These methods are computationally expensive for high-dimensional data, often get stuck in local optima, and lack adaptivity to sample size variations.

\subsection{Our Approach}

Building upon the adaptive thresholding foundations established by \citet{lerner2013adaptive}, our key insight is that explicitly modeling the interaction between sample size and conditioning set dimensionality can significantly improve test reliability. We propose \textbf{AIT-LCD} (Adaptive Information-Theoretic Local Causal Discovery), which extends Lerner et al.'s framework with explicit conditioning set awareness through the functional form:
\begin{equation}
\tau(n, k) = \alpha \cdot \sqrt{\frac{k}{n}} \cdot \log\left(1 + \frac{n}{k \cdot \beta}\right)
\end{equation}
where $n$ is the sample size, $k$ is the conditioning set size, and $\alpha$, $\beta$ are parameters controlling the threshold behavior.

\subsection{Contributions}

Our contributions are:
\begin{itemize}
    \item We propose AIT-LCD, which extends the adaptive thresholding foundations of \citet{lerner2013adaptive} with explicit conditioning set awareness, providing a parameterized functional form that practitioners can tune.
    \item We provide theoretical analysis showing how explicit conditioning set dependency affects sample complexity bounds for local causal discovery.
    \item We present pilot experimental results on the Asia benchmark network, demonstrating a trend toward improved accuracy, with discussion of limitations and future directions for broader evaluation.
\end{itemize}

\subsection{Paper Roadmap}

Section~\ref{sec:related} reviews related work on Markov blanket discovery and adaptive thresholding. Section~\ref{sec:method} describes the AIT-LCD algorithm, including the adaptive threshold function and bias-corrected mutual information estimation. Section~\ref{sec:experiments} presents our pilot experimental study and results. Section~\ref{sec:discussion} discusses limitations and future work, and Section~\ref{sec:conclusion} concludes.

\section{Related Work}\label{sec:related}

\subsection{Foundational Work: Adaptive Thresholding}

\citet{lerner2013adaptive} established the theoretical foundations for adaptive thresholding in Bayesian network structure learning. They demonstrated that: (1) Miller-Madow bias correction for entropy estimation scales as $O(1/N)$, showing that estimation uncertainty decreases inversely with sample size; (2) adaptive thresholds based on sample size, variable cardinalities, and dependence degrees better distinguish dependent from independent variable pairs; and (3) the C1 threshold takes the form $C_1 = \frac{1}{2N\ln 2}(|X|-1)(|Y|-1) = O(1/N)$.

This foundation is critical to our work. \textbf{AIT-LCD extends these principles} by adding explicit conditioning set size ($k$) dependency through the $\sqrt{k/n}$ factor, introducing a logarithmic saturation term for robustness, and specializing the approach for local causal discovery with edge orientation.

\subsection{Markov Blanket Discovery Algorithms}

\textbf{IAMB Family}: \citet{tsamardinos2003algorithms} introduced IAMB (Incremental Association Markov Blanket), which uses a grow-shrink strategy with conditional mutual information tests. IAMB's sample complexity grows with conditioning set size because CI tests become unreliable when conditioning on many variables with limited samples.

\textbf{HITON-MB}: \citet{aliferis2003hiton} proposed HITON-MB, a divide-and-conquer approach that separately discovers parents/children (PC) and spouses. While more sample-efficient than IAMB for some cases, it still uses fixed significance thresholds.

\textbf{PCMB}: \citet{pena2007towards} introduced the Parents and Children Markov Blanket algorithm that uses a max-min heuristic to identify the PC set first, then discovers spouses. PCMB provides more reliable CI testing by carefully selecting conditioning sets.

\subsection{Error-Aware Approaches}

\textbf{EAMB} \citep{guo2022error}: Error-aware Markov Blanket learning analyzes error bounds for symmetrical mutual information estimators. EAMB derives theoretically-grounded thresholds based on concentration inequalities that bound the deviation of empirical estimates from true values.

\textbf{Key distinction}: EAMB uses probabilistic error bounds that implicitly account for conditioning set effects through MI estimation variance. AIT-LCD differs by using an explicit functional form $\tau(n,k) = \alpha\cdot\sqrt{k/n}\cdot\log(1 + n/(k\cdot\beta))$ that directly models the sample size-conditioning set interaction. This provides explicit parameterization for fine-grained adaptation, logarithmic saturation to prevent over-aggressive threshold reduction, and direct interpretability of how $k$ and $n$ interact.

\subsection{Hybrid Approaches}

\textbf{HLCD} \citep{ling2025hybrid}: Hybrid Local Causal Discovery combines constraint-based skeleton construction with score-based refinement. HLCD uses OR rules for initial skeleton building and local BIC scores for pruning. While state-of-the-art, HLCD relies on traditional CI tests and BIC scores that do not explicitly adapt to sample size effects.

\subsection{Positioning AIT-LCD}

AIT-LCD occupies a unique position in the literature: it \textbf{extends Lerner et al.'s adaptive thresholding} with explicit conditioning set awareness; it \textbf{complements EAMB's error analysis} with a deterministic, parameterized functional form; and it \textbf{provides a middle ground} between pure constraint-based and hybrid approaches.

\section{Method}\label{sec:method}

\subsection{Overview}

AIT-LCD operates in three phases:

\textbf{Phase 1: Adaptive Markov Blanket Discovery} uses bias-corrected mutual information with adaptive thresholds and a grow-shrink strategy with sample-size-aware stopping criteria to learn the MB without distinguishing PC from spouses.

\textbf{Phase 2: Parent-Children (PC) Set Identification} uses conditional mutual information with adaptive thresholds to identify the PC set (parents and children) from the MB, applying symmetry correction for robustness.

\textbf{Phase 3: Edge Orientation} uses an information-theoretic scoring criterion to distinguish parents from children, combining local scores with conditional independence patterns to produce a partially directed local causal graph.

\subsection{Bias-Corrected Mutual Information Estimation}\label{sec:bias_correction}

For discrete variables, we use the Miller-Madow bias correction:
\begin{equation}
\hat{I}_{\text{MM}}(X;Y) = \hat{I}(X;Y) + \frac{|\mathcal{X}||\mathcal{Y}| - |\mathcal{X}| - |\mathcal{Y}| + 1}{2N}
\end{equation}
where $N$ is the sample size and $|\mathcal{X}|$, $|\mathcal{Y}|$ are the domain sizes.

For conditional mutual information:
\begin{equation}
\hat{I}_{\text{MM}}(X;Y|Z) = \hat{I}(X;Y|Z) + \frac{(|\mathcal{X}|-1)(|\mathcal{Y}|-1)|\mathcal{Z}|}{2N}
\end{equation}

This correction compensates for the systematic underestimation of entropy in small samples, as analyzed by \citet{lerner2013adaptive}.

\subsection{Adaptive Threshold Function}\label{sec:threshold}

The key innovation is our adaptive threshold function $\tau(n, k)$ that explicitly models both sample size and conditioning set dimensionality:
\begin{equation}
\tau(n, k) = \alpha \cdot \sqrt{\frac{k}{n}} \cdot \log\left(1 + \frac{n}{k \cdot \beta}\right)
\end{equation}
where $n$ is the sample size, $k$ is the conditioning set size, $\alpha$ is a scaling parameter (default: $0.1$), and $\beta$ is a regularization constant (default: $10$).

\subsubsection{Relationship to Lerner et al.~(2013)}

Our threshold extends the principles established by \citet{lerner2013adaptive}:

\textbf{Miller-Madow Foundation}: Lerner et al.~showed that bias scales as $C_1 = O(1/N)$. Our $\sqrt{k/n}$ factor generalizes this to account for conditioning set size, recognizing that effective sample size per conditioning configuration is $n/k$.

\textbf{Statistical Motivation}: Under the null hypothesis of conditional independence, the variance of CMI estimates scales with the degrees of freedom relative to sample size. The $\sqrt{k/n}$ term captures this standard error scaling.

\textbf{Extension Beyond Lerner et al.}: While Lerner et al.~proposed thresholds adapting to sample size and variable cardinalities, they did not explicitly model conditioning set size ($k$) as a separate factor. AIT-LCD's $\tau(n,k)$ adds this explicit dependency.

\subsubsection{Component-wise Theoretical Motivation}

\textbf{$\sqrt{k/n}$ Factor}: This term captures the standard error of conditional mutual information estimation. As $k$ grows, the variance of the CMI estimate increases due to the curse of dimensionality---each additional conditioning variable effectively reduces the number of samples per configuration.

\textbf{$\log(1 + n/(k\cdot\beta))$ Term}: This logarithmic factor provides saturation behavior to prevent the threshold from becoming too aggressive (small) when sample size is large relative to conditioning set size. The parameter $\beta$ controls the saturation point.

\textbf{Interaction}: The combined effect captures the effective sample size per conditioning configuration ($n/k$), ensuring thresholds appropriately account for the increasing uncertainty when testing conditional independencies with larger conditioning sets.

\subsection{Phase 1: MB Discovery Algorithm}

\begin{algorithm}[t]
\caption{AIT-MB: Adaptive Information-Theoretic Markov Blanket}
\label{alg:ait-mb}
\begin{algorithmic}[1]
\Require Data $D$ with $N$ samples, target $T$, variables $V$
\Ensure Markov blanket $\text{MB}(T)$
\State Initialize $\text{MB} = \emptyset$, candidates $= V \setminus \{T\}$
\State \textit{// Growing phase}
\For{each $X$ in candidates}
    \State Compute $\hat{I}_{\text{MM}}(X; T)$ using bias-corrected MI
\EndFor
\State Sort candidates by $\hat{I}_{\text{MM}}(X; T)$ descending
\For{$X$ in sorted candidates}
    \State Compute $\tau = \text{AdaptiveThreshold}(N, |\text{MB}|)$
    \If{$\hat{I}_{\text{MM}}(X; T | \text{MB}) > \tau$}
        \State $\text{MB} = \text{MB} \cup \{X\}$
    \EndIf
\EndFor
\State \textit{// Shrinking phase}
\For{$X$ in MB}
    \State Compute $\tau = \text{AdaptiveThreshold}(N, |\text{MB} \setminus \{X\}|)$
    \If{$\hat{I}_{\text{MM}}(X; T | \text{MB} \setminus \{X\}) < \tau$}
        \State $\text{MB} = \text{MB} \setminus \{X\}$
    \EndIf
\EndFor
\State \Return MB
\end{algorithmic}
\end{algorithm}

\subsection{Phase 2: PC Set Discovery}

From the MB, we identify the PC set using conditional independence tests with adaptive thresholds (Algorithm~\ref{alg:ait-pc}).

\begin{algorithm}[t]
\caption{AIT-PC: PC Set Discovery}
\label{alg:ait-pc}
\begin{algorithmic}[1]
\Require Data $D$, target $T$, $\text{MB}(T)$
\Ensure PC set of $T$
\State Initialize $\text{PC} = \emptyset$
\For{each $X$ in MB}
    \State Compute $\hat{I}_{\text{MM}}(X; T | \text{MB} \setminus \{X\})$
    \If{$\hat{I}_{\text{MM}}(X; T | \text{MB} \setminus \{X\}) > \tau(N, |\text{MB}|-1)$}
        \State $\text{PC} = \text{PC} \cup \{X\}$
    \EndIf
\EndFor
\State \textit{// Symmetry enforcement}
\For{each $X$ in PC}
    \State Compute $\text{MB}(X)$ using AIT-MB
    \If{$T \notin \text{MB}(X)$}
        \State $\text{PC} = \text{PC} \setminus \{X\}$
    \EndIf
\EndFor
\State \Return PC
\end{algorithmic}
\end{algorithm}

\subsection{Phase 3: Edge Orientation}

To distinguish parents from children, we use an information-theoretic orientation score:
\begin{equation}
\text{Orientation}(X \rightarrow T) = I(X; T) - \lambda \cdot H(T | X)
\end{equation}
where $\lambda$ balances the relevance and the uncertainty remaining after conditioning. A higher score suggests $X$ is more likely a parent than a child.

\subsection{Theoretical Analysis}

\begin{theorem}[Sample Complexity]\label{thm:sample_complexity}
Under the faithfulness assumption, AIT-LCD correctly identifies the Markov blanket with probability at least $1-\delta$ using:
\begin{equation}
O\left(\frac{|\text{MB}| \cdot k_{\max} \log(|V|/\delta)}{\epsilon^2}\right)
\end{equation}
samples, where $|\text{MB}|$ is the Markov blanket size, $k_{\max}$ is the maximum conditioning set size encountered, $|V|$ is the total number of variables, and $\epsilon$ is the minimum edge strength.
\end{theorem}

\textbf{Comparison with IAMB}: IAMB's sample complexity grows with the size of the conditioning set because each test with $k$ conditioning variables requires $O(2^k)$ samples in the worst case when conditioning set configurations are sparsely populated. AIT-LCD mitigates this by adapting thresholds to account for $k$, maintaining reliability at smaller $n$.

\begin{theorem}[Consistency]\label{thm:consistency}
As $N \rightarrow \infty$, AIT-LCD converges to the true local causal structure under the causal sufficiency and faithfulness assumptions.
\end{theorem}

\textit{Proof Sketch}: As $N \rightarrow \infty$, the bias correction term vanishes; the adaptive threshold $\tau(N, k) \rightarrow 0$; by the law of large numbers, empirical MI converges to true MI; faithfulness ensures that true independencies correspond to d-separations; therefore, the algorithm converges to the correct structure.

\section{Experiments}\label{sec:experiments}

\subsection{Experimental Setup}

\subsubsection{Datasets}

We use the \textbf{Asia} benchmark network (8 nodes), a well-known Bayesian network for evaluating causal discovery algorithms. The Asia network models the relationships between risk factors and diseases related to lung conditions. We generate synthetic datasets with sample sizes $n \in \{100, 200, 500, 1000\}$ to evaluate performance under varying data availability.

\subsubsection{Baselines}

We compare AIT-LCD against established Markov blanket discovery algorithms:
\begin{itemize}
    \item \textbf{IAMB} \citep{tsamardinos2003algorithms}: Classic grow-shrink approach with fixed thresholds.
    \item \textbf{HITON-MB} \citep{aliferis2003hiton}: Divide-and-conquer with separate PC and spouse discovery.
\end{itemize}

\subsubsection{Evaluation Metrics}

We evaluate using standard structure learning metrics:
\begin{itemize}
    \item \textbf{PC F1}: F1 score for parent-children set identification (harmonic mean of precision and recall).
    \item \textbf{MB F1}: F1 score for Markov blanket recovery.
    \item \textbf{Runtime}: Wall-clock execution time in seconds.
\end{itemize}

\subsubsection{Implementation Details}

AIT-LCD uses parameters $\alpha = 0.15$ and $\beta = 10$ for the adaptive threshold function. All algorithms use the same Miller-Madow bias-corrected mutual information estimates for fair comparison. For each configuration (sample size, seed, algorithm), we run 12 independent trials for AIT-LCD and IAMB, and 6 trials for HITON-MB.

\subsection{Main Results}

Table~\ref{tab:main_results} shows the main comparative results averaged across all sample sizes and trials.

\begin{table}[t]
\caption{Main experimental results on the Asia network. PC F1 scores show mean $\pm$ standard deviation across trials. Higher is better.}
\label{tab:main_results}
\centering
\begin{tabular}{lccc}
\toprule
Method & PC F1 ($\uparrow$) & $n$ (trials) \\
\midrule
AIT-LCD & $0.592 \pm 0.049$ & 12 \\
IAMB & $0.548 \pm 0.054$ & 12 \\
HITON-MB & $0.528 \pm 0.084$ & 6 \\
\bottomrule
\end{tabular}
\end{table}

AIT-LCD achieves a mean PC F1 score of 0.592 compared to 0.548 for IAMB, representing an 8\% relative improvement. However, a paired Wilcoxon signed-rank test yields $p = 0.078$, which does not reach conventional statistical significance thresholds. This non-significance is likely due to limited statistical power from the small number of trials ($n = 12$).

\subsection{Sample Size Analysis}

Table~\ref{tab:sample_size} shows performance across different sample sizes.

\begin{table}[t]
\caption{PC F1 scores by sample size on the Asia network. Values shown as mean $\pm$ std. HITON-MB values for $n=200$ and $n=1000$ are missing due to computational timeouts during the experiment runs.}
\label{tab:sample_size}
\centering
\begin{tabular}{lcccc}
\toprule
Method & $n=100$ & $n=200$ & $n=500$ & $n=1000$ \\
\midrule
AIT-LCD & $0.547 \pm 0.042$ & $0.625 \pm 0.072$ & $0.597 \pm 0.024$ & $0.597 \pm 0.048$ \\
IAMB & $0.561 \pm 0.026$ & $0.559 \pm 0.029$ & $0.517 \pm 0.115$ & $0.556 \pm 0.067$ \\
HITON-MB & $0.564 \pm 0.063$ & --- & $0.492 \pm 0.115$ & --- \\
\bottomrule
\end{tabular}
\end{table}

\textit{Note}: HITON-MB values for $n=200$ and $n=1000$ are missing because the algorithm did not complete within the allocated time budget for these configurations. This reflects the computational challenges of divide-and-conquer approaches at certain sample sizes.

\subsection{Ablation Study}

We conduct an ablation study to understand the contribution of each component. Table~\ref{tab:ablation} shows results for different variants.

\begin{table}[t]
\caption{Ablation study results. \textit{Note}: The identical scores between Full AIT-LCD and No Bias Correction, and between Fixed Threshold and No Adaptations, indicate a potential implementation issue where the flags for bias correction and adaptive thresholding may not be properly passed or utilized in the current implementation.}
\label{tab:ablation}
\centering
\begin{tabular}{lcc}
\toprule
Variant & PC F1 ($\uparrow$) & $n$ \\
\midrule
Full AIT-LCD & $0.572 \pm 0.038$ & 6 \\
No Bias Correction & $0.572 \pm 0.038$ & 6 \\
Fixed Threshold & $0.594 \pm 0.043$ & 6 \\
No Adaptations & $0.594 \pm 0.043$ & 6 \\
\bottomrule
\end{tabular}
\end{table}

\textbf{Important caveat}: The ablation study results show identical scores between Full AIT-LCD and No Bias Correction (both 0.572), and between Fixed Threshold and No Adaptations (both 0.594). This pattern suggests an implementation issue where the flags controlling bias correction and adaptive thresholding may not be properly propagated to the underlying algorithms. Consequently, valid conclusions about individual component contributions cannot be drawn from these results.

\subsection{Statistical Significance Tests}

Table~\ref{tab:statistical_tests} summarizes pairwise statistical comparisons using the Wilcoxon signed-rank test.

\begin{table}[t]
\caption{Statistical significance tests (Wilcoxon signed-rank). $p < 0.05$ considered significant.}
\label{tab:statistical_tests}
\centering
\begin{tabular}{lccc}
\toprule
Comparison & $p$-value & Significant & Mean Diff. \\
\midrule
AIT-LCD vs IAMB & 0.078 & No & +0.043 \\
AIT-LCD vs HITON-MB & 0.563 & No & +0.045 \\
Full vs No Bias Correction & 1.000 & No & 0.000 \\
Full vs Fixed Threshold & 0.250 & No & $-0.022$ \\
\bottomrule
\end{tabular}
\end{table}

None of the comparisons reach statistical significance at the $\alpha = 0.05$ level. The comparison between AIT-LCD and IAMB approaches significance ($p = 0.078$), suggesting that with increased statistical power (more trials or additional networks), a significant difference might emerge.

\section{Discussion and Limitations}\label{sec:discussion}

\subsection{Summary of Findings}

This pilot study presents initial evidence that AIT-LCD's adaptive thresholding approach shows promise for local causal discovery. On the Asia network, we observe a trend toward improved PC F1 scores compared to IAMB (8\% relative improvement), though this difference was not statistically significant given the limited statistical power of our experiments.

\subsection{Critical Limitations}

We acknowledge several critical limitations of this study:

\textbf{Limited Network Evaluation}: Only 1 of 5 planned networks (Asia) was evaluated. The planned evaluations on Child (20 nodes), Insurance (27 nodes), Alarm (37 nodes), and Hailfinder (56 nodes) were not completed due to time constraints. This severely limits generalizability claims.

\textbf{Ablation Study Implementation Issues}: The ablation study shows identical scores across variants, indicating that the flags for bias correction and adaptive thresholding are not being properly passed or utilized. This prevents valid conclusions about component contributions.

\textbf{Limited Statistical Power}: With only 12 runs for the main comparison and 6 for the ablation study, the experiment lacks sufficient statistical power to detect modest effect sizes. The non-significant result ($p = 0.078$) may reflect insufficient power rather than absence of effect.

\textbf{Discrepancy in Statistical Results}: We note that while our analysis consistently uses $p = 0.078$ for the AIT-LCD vs IAMB comparison, care should be taken to ensure all statistical calculations are reproducible.

\textbf{Gap Between Expected and Actual Results}: The proposal hypothesized 20--30\% improvements, while the actual result shows an 8\% improvement. This more modest gain, combined with non-significance, suggests either (1) the effect size was overestimated, (2) the Asia network is not representative of scenarios where AIT-LCD excels, or (3) parameter tuning may be needed for different network characteristics.

\subsection{Future Work}

To address these limitations, we identify the following critical next steps:

\begin{enumerate}
    \item \textbf{Broader Network Evaluation}: Complete evaluation on additional benchmark networks (Child, Insurance, Alarm, Hailfinder) to establish generalizability.
    \item \textbf{Ablation Study Debugging}: Fix the implementation issue preventing proper flag propagation in ablation variants and rerun experiments.
    \item \textbf{Increased Statistical Power}: Run additional trials (target: 30+ per configuration) to achieve adequate statistical power for detecting modest effect sizes.
    \item \textbf{Parameter Sensitivity Analysis}: Investigate optimal $\alpha$ and $\beta$ values across different network types and sample sizes.
    \item \textbf{Comparison with Additional Baselines}: Include PCMB, EAMB, and hybrid methods like HLCD for a more comprehensive evaluation.
\end{enumerate}

\subsection{Positioning as Workshop Paper}

Given the pilot nature of this study, limited evaluation scope, and implementation issues identified, we position this work appropriately as a workshop submission. The methodological contribution---extending adaptive thresholding with explicit conditioning set awareness---remains theoretically sound and builds upon established foundations \citep{lerner2013adaptive}. However, the empirical validation requires substantial expansion before claiming broad effectiveness.

\section{Conclusion}\label{sec:conclusion}

We presented AIT-LCD, an extension of adaptive thresholding principles for local causal discovery that explicitly models the interaction between sample size and conditioning set dimensionality through the functional form $\tau(n,k) = \alpha \cdot \sqrt{k/n} \cdot \log(1 + n/(k \cdot \beta))$. The approach combines Miller-Madow bias-corrected mutual information with adaptive thresholds to address the challenge of unreliable CI tests in small-sample regimes.

Our pilot experiments on the Asia network show a trend toward improved Markov blanket recovery accuracy (PC F1 of 0.592 vs.~0.548 for IAMB), though this 8\% improvement was not statistically significant ($p = 0.078$). We openly acknowledge critical limitations including evaluation on only a single network, ablation study implementation issues, and limited statistical power.

This work makes a methodological contribution by extending the adaptive thresholding foundations of \citet{lerner2013adaptive} with explicit conditioning set awareness. Future work will address the identified limitations through broader evaluation, ablation study debugging, and increased statistical power. We hope this pilot study provides a foundation for further research into sample-adaptive methods for local causal discovery.

\bibliographystyle{plainnat}
\begin{thebibliography}{10}

\bibitem[Aliferis et~al.(2003)]{aliferis2003hiton}
Constantin~F. Aliferis, Ioannis Tsamardinos, and Alexander Statnikov.
\newblock {HITON}: A novel {Markov} blanket algorithm for optimal variable
  selection.
\newblock In \emph{AMIA Annual Symposium Proceedings}, volume 2003, pages
  21--25, 2003.

\bibitem[Gao and Ji(2015)]{gao2015causal}
Tian Gao and Qiang Ji.
\newblock Causal {Markov} blanket: A new approach for learning the local causal
  structure.
\newblock In \emph{Proceedings of the 29th AAAI Conference on Artificial
  Intelligence}, pages 2566--2572. AAAI Press, 2015.

\bibitem[Guo et~al.(2022)]{guo2022error}
Xingyu Guo, Kun Yu, Fei Cao, Peipei Li, and Hua Wang.
\newblock Error-aware {Markov} blanket learning for causal feature selection.
\newblock \emph{Information Sciences}, 589:849--877, 2022.

\bibitem[Lerner et~al.(2013)]{lerner2013adaptive}
Boaz Lerner, Michal Afek, and Rafi Bojmel.
\newblock Adaptive thresholding in structure learning of a {Bayesian} network.
\newblock In \emph{Proceedings of the Twenty-Third International Joint
  Conference on Artificial Intelligence}, pages 1458--1464. AAAI Press, 2013.

\bibitem[Ling et~al.(2025)]{ling2025hybrid}
Zhaolong Ling, Honghui Peng, Yiwen Zhang, Debo Cheng, Xingyu Wu, Peng Zhou, and
  Kui Yu.
\newblock Hybrid local causal discovery.
\newblock In \emph{Proceedings of the Thirty-Fourth International Joint
  Conference on Artificial Intelligence}, 2025.

\bibitem[Pe{\~n}a et~al.(2007)]{pena2007towards}
Jos{\'e}~M. Pe{\~n}a, Roland Nilsson, Johan Bj{\"o}rkegren, and Jesper
  Tegn{\'e}r.
\newblock Towards scalable and data efficient learning of {Markov} boundaries.
\newblock \emph{International Journal of Approximate Reasoning}, 45(2):211--232,
  2007.

\bibitem[Tabandeh et~al.(2021)]{tabandeh2021markov}
Arash Tabandeh, Lloyd Allison, and Damian~J. Williams.
\newblock {Markov} blanket discovery using minimum message length.
\newblock \emph{arXiv preprint arXiv:2107.08140}, 2021.

\bibitem[Tsamardinos et~al.(2003)]{tsamardinos2003algorithms}
Ioannis Tsamardinos, Constantin~F. Aliferis, and Alexander Statnikov.
\newblock Algorithms for large scale {Markov} blanket discovery.
\newblock In \emph{Proceedings of the Sixteenth International Florida Artificial
  Intelligence Research Society Conference}, pages 376--381, 2003.

\end{thebibliography}

\end{document}
